% Generated draft: requires mathematical and source review.
\documentclass[11pt]{article}
\usepackage[margin=25mm]{geometry}
\usepackage{fontspec}
\setmainfont{Latin Modern Roman}
\usepackage{amsmath,amssymb,mathtools,unicode-math}
\setmathfont{Latin Modern Math}
\usepackage{enumitem,needspace,xurl,hyperref}
\setlength{\parindent}{0pt}
\setlength{\parskip}{5pt}
\setlength{\emergencystretch}{3em}
\newcommand{\N}{\mathbb{N}}
\newcommand{\Z}{\mathbb{Z}}
\newcommand{\Q}{\mathbb{Q}}
\newcommand{\R}{\mathbb{R}}
\newcommand{\C}{\mathbb{C}}
\newcommand{\F}{\mathbb{F}}
\newcommand{\A}{\mathbb{A}}
\newcommand{\PP}{\mathbb{P}}
\newcommand{\E}{\mathbb{E}}
\newcommand{\eps}{\varepsilon}
\newcommand{\dd}{\,\mathrm{d}}
\begin{document}
\section*{UnsolvedMath ID: 30005607}
\textbf{Source problem code:} catalogue entry
\subsection{Definitions and mathematical statement}
Let \(S^2=\{x\in\R^3:|x|=1\}\) carry its standard round metric, with area measure \(\dd A\) and total area \(4\pi\). A domain is a nonempty connected open subset of \(S^2\). For \(0<A<4\pi\), let \(C_A\) be a spherical cap of area \(A\); its center is immaterial by rotational invariance.

For a domain \(\Omega\), define the variational quantity
\[
\mu_2(\Omega)=\inf_{\substack{u\in H^1(\Omega)\setminus\{0\}\\ \int_\Omega u\,\dd A=0}}
\frac{\int_\Omega |\nabla_{S^2}u|^2\,\dd A}{\int_\Omega |u|^2\,\dd A}.
\]
When the Neumann Laplacian has discrete spectrum, this is its second eigenvalue under the convention
\[
0=\mu_1(\Omega)\leq\mu_2(\Omega)\leq\mu_3(\Omega)\leq\cdots,
\]
with eigenvalues counted with multiplicity. On a smooth boundary the Neumann condition is \(\partial_\nu u=0\), where \(\nu\) is the outward unit conormal, and the operator is \(-\Delta_{S^2}\).

\textbf{Conjecture 1, as supplied, with unresolved regularity qualifications.} For every prescribed area \(A\in(0,4\pi)\) and every admissible simply connected domain \(\Omega\subset S^2\) of area \(A\),
\[
\mu_2(\Omega)\leq\mu_2(C_A).
\]
Thus a spherical cap is a maximizer; uniqueness of maximizers is not asserted.

The supplied record does not specify the admissible boundary regularity or spectral hypotheses. The variational definition above supplies one precise interpretation, but its intended use for arbitrary irregular domains is unconfirmed; without suitable compactness, the infimum need not be an eigenvalue. The displayed area range treats proper caps. The source's phrase ``up to 100 percent'' does not specify whether the full-area endpoint \(A=4\pi\), the whole sphere, or full-area proper domains are included.

The source describes this as improving a preceding theorem's area fraction from \(0.94\) to \(1.00\), but does not supply that theorem's complete hypotheses. Its proposed approach is to retain qualitative information discarded in the existing proof; that information is not identified. The extracted text also includes a figure caption reporting numerical counterexamples attributed to Martinet, reference [20], at area fractions \(0.64\) and \(0.86\) of \(S^2\). Their domain class and the inequality being contradicted are not supplied, so they cannot be identified here as counterexamples to the simply connected conjecture.
\subsection{Short English statement}
For each prescribed area strictly between zero and the area of the round sphere, the conjecture says that a spherical cap maximizes the second Neumann eigenvalue among simply connected domains of that area, counting the constant eigenfunction first. Boundary regularity, the full-area endpoint, and the context of the extracted numerical counterexamples require clarification. The supplied open status is dated curation, not an independently verified conclusion.
\subsection{Sources}
\begin{itemize}
\item Local UnsolvedMath registry: permanent ID \(30005607\), source problem code \texttt{catalogue entry}, title ``Unrestricted Spherical-Cap Maximization of the Second Neumann Eigenvalue.'' The supplied record gives proposed year 2023 and no named proposer.
\item Supplied report citation: \emph{Geometric Spectral Theory (2024)}, Oberwolfach Reports, pp.\ 2017--2106, DOI \texttt{10.4171/owr/2023/36}. \url{https://doi.org/10.4171/owr/2023/36}. The original excerpt labels the assertion Conjecture 1 and includes a page marker for p.\ 2066 of Oberwolfach Report 36/2023. The citation's 2024 date and the report's 2023 designation are retained without reconciliation.
\item Supplied literature source: ``Oberwolfach Report 36.'' \url{https://doi.org/10.4171/owr/2023/36}.
\item Supplied literature source: ``Bucur et al. counterexample and open problems.'' \url{https://arxiv.org/abs/2503.15385}. The supplied assessment describes this as a 2025 counterexample paper and says it explicitly leaves the simply connected range above 94 percent unresolved; this characterization has not been independently verified.
\item Literature assessment dated exactly \texttt{2026-08-22}: the supplied curation labels the problem open, states that the 94-percent theorem is known, and states that extending the result to all simply connected domains remains open. This is supplied dated curation, not an independently verified literature assessment. The record itself requests independent verification of time-sensitive status and cited claims.
\item The registry labels the statement \texttt{corrected\_verified} and reports reconstruction and checking against the official Oberwolfach report. These are supplied verification claims, not verification performed for this draft. The extracted caption attributes numerical examples to Martinet [20], whose full bibliographic reference is absent.
\end{itemize}
\end{document}
